\documentclass[11pt]{letter}
\usepackage[T1]{fontenc}
\usepackage{mathptmx}
\usepackage[a4paper,margin=2.5cm]{geometry}
\signature{Pedro S. Romano\\on behalf of all authors}
\address{Pedro S. Romano\\Iscte -- Instituto Universit\'ario de Lisboa\\1649-026 Lisbon, Portugal\\Pedro\_Romano@iscte-iul.pt}
\begin{document}
\begin{letter}{The Editor-in-Chief\\IEEE Access}
\opening{Dear Editor,}

We submit the manuscript \emph{Configuring LLM-Based Grading in Computer Science: An Empirical Study of
Reasoning, Guidance, Scope, and Cost} for consideration as a regular article in IEEE Access.

Instructors and institutions adopting large language models to grade student work face practical choices:
whether the model should reason before scoring, what guidance it should receive, how much of an exam it
should see at once, and which model is worth its cost. Prior studies have examined these choices one at a
time. We compared them under one protocol, with five models on four datasets of short answers and code,
including exams in Portuguese from a course two of us taught, and measured agreement with teachers, consistency across
repeated runs, and cost together.

The cheap changes proved the effective ones: giving the model a rubric or a reference answer, and asking it
for partial credit. Reasoning rarely justified its cost; its large gains came on code, for models that score
many partially correct answers as zero, a failure we describe and measure. From these results we built a
simple configuration guide, tested it against every other configuration in the study, and validated it on
two datasets that played no part in building it.

We believe the paper suits the broad engineering readership of IEEE Access: it gives practitioners who build
or deploy grading systems a tested configuration and its cost. The supplementary material includes every
prompt template, the values behind every figure, and a package with the models' grades on the public datasets
and a script that recomputes the reported results.

The manuscript is original, has not been published, and is not under consideration elsewhere. All authors
have approved it, and they declare no competing interests. The manuscript includes the statement on research with human subjects and the statement
on the use of AI tools required by IEEE.

Thank you for considering our submission.

\closing{Sincerely,}
\end{letter}
\end{document}
